\chapter{Alternative-Data Pipelines}\label{ch:ml:alternative-data-pipelines}

A card-spending panel nowcasts a hundred retailers' year-on-year sales growth within 3.3 percentage points on average for
three years; in the fourth it misses by 22.5. Nothing changed at the retailers. The vendor had changed its bank partner,
and the new partner's cardholders were older and lived elsewhere. Book~7 (chapter~12) taught how to decide whether to buy
a dataset; Book~8 (chapter~16) traded one. This chapter is the plumbing in between, where most of the work and most of the
losses are: ingesting deliveries that are sometimes wrong, turning merchant strings into securities, correcting a panel
that is not the population, and telling a history the vendor lived through from one it reconstructed. Every number comes
from a synthetic panel whose truth is known, built so that each stage has something to catch.

\section{Ingestion and validation}

\begin{definition}[Ingestion pipeline, schema validation]\label{def:ml:alternative-data-pipelines:ingest}
An \emph{ingestion pipeline}\index{ingestion pipeline} is the code that receives a vendor's deliveries, checks them, maps
them to the firm's identifiers and stores them with the time each value became known. \emph{Schema
validation}\index{schema validation} checks each record against a declared schema: every field present, of its type and in
its range, and no record duplicated; records that fail are rejected with a reason, never silently repaired.
\end{definition}

The chapter's world has 100 retailers over 24 quarters. Their true sales are spent by a population divided into six cells,
three age bands by two regions, with known shares; each retailer draws its customers from the cells in its own proportions,
and spending in each cell grows at its own rate. The vendor's panel holds 50\,000 cardholders from one bank, tilted young
(45\% in the youngest band against 30\% of the population) and drifting younger by a point a quarter; at quarter 16 the
vendor changes partner and the panel becomes 30\,000 cardholders, half in the oldest band and three-quarters in one region.
The vendor launches at quarter 12, delivering quarters 0 to 11 at once, then one quarter's data a quarter after it ends.

Each delivery is rows of quarter, merchant string, age band, region, spend and panel count, one retailer's spend in a cell
split across two or three store-level strings: 36\,766 rows in all, with planted defects. \omterm{def:ml:alternative-data-pipelines:ingest}{Schema validation}
(\cref{lst:ml:alt:validate}) rejects 121 rows with a missing field, 69 with a negative spend, 64 with the panel count
delivered as text, and 111 duplicates, every planted defect and nothing else. A row can be valid and the delivery wrong: the
delivery for quarter 19 arrives in cents. A second check compares each delivery's total with the previous one and holds it
for review when the ratio leaves $[0.5, 2]$; it holds that delivery and no other, and review rescales it. Checks on the
delivery, not only on the rows, are what catch a vendor's change of units, a lost file or a file sent twice.

\section{Mapping to securities}

\begin{definition}[Entity resolution, record linkage]\label{def:ml:alternative-data-pipelines:er}
\emph{Record linkage}\index{record linkage} decides which records in two files describe the same entity when no common
identifier exists, by comparing their attributes (Fellegi and Sunter, 1969). \emph{Entity resolution}\index{entity
resolution} is the task applied to a dataset's entity strings (merchants, issuers, apps, vessels) against a reference list,
here the security master (Book~7, chapter~4): normalise, compare, and decide to match, to send to review, or to leave
unmatched.
\end{definition}

A card network prints ``SQ *BOREALIS RET'', ``BOREALIS RETAIL \#4417'', ``BOREALI RETA STR 313'' or ``BOREALIS 5120'', and
the model needs the permanent identifier of Borealis Retail, not a string. The resolver of \cref{lst:ml:alt:resolve}
strips processor prefixes, store numbers and filler words, then scores each alias by its tokens, counting a token found as a
prefix of three letters or more (card networks truncate) as a full match. Of the 600 strings of the retailers, exact matching
on the normalised name finds 300; the resolver matches 500 automatically, all correctly, and sends the other 100 to review:
the strings that give only the name's first word, which up to three retailers share (Borealis Retail, Borealis Stores,
Borealis Outfitters). Of ten strings from unrelated merchants (``DELTA AIR LINES'', ``SUMMIT DENTAL CARE''), none is matched,
six go to review and four are rejected outright. The review queue is where a person earns their place in the pipeline:
matching a store to the wrong retailer contaminates its sales silently, while a string left for review only delays it.

\section{Panel bias and reweighting}

\begin{definition}[Panel bias, panel reweighting]\label{def:ml:alternative-data-pipelines:panel}
\emph{Panel bias}\index{panel bias} is the difference between a statistic computed on a data vendor's panel (cardholders,
devices, users) and the same statistic on the population it stands for, caused by who is in the panel. \emph{Panel
reweighting}\index{panel reweighting} gives each panel member a weight so that the weighted panel matches known population
totals; \emph{raking} (iterative proportional fitting) does it with only the margins, adjusting the weights to each margin
in turn until all match (Deming and Stephan, 1940).
\end{definition}

The raw nowcast divides each retailer's panel spend by the number of panelists and reads year-on-year growth from it. While
the panel drifts slowly, the error is modest (3.3 points on average over quarters 4 to 15): a retailer whose young customers
spend more each year looks, in a panel getting younger, as if it grows faster than it does. In the four quarters after the
partner change, each year-on-year comparison sets the new panel against the old, and the error jumps to 22.5 points
(\cref{fig:ml:alt:errors}); it returns to 2.7 once both sides of the comparison come from the new panel. Raking the six cells
to the population's age and region margins (\cref{lst:ml:alt:rake}) removes the composition effect: 1.7 points before the
change, 2.7 in the change year, 2.6 after. The change year's residual comes from the new panel's small cells (1\,125
cardholders in the thinnest), whose sampling noise no reweighting removes.

\begin{omfigure}
\begin{tikzpicture}
\begin{axis}[width=11.5cm,height=6cm,xlabel={quarter},ylabel={mean absolute error (points)},
  tick label style={font=\small},label style={font=\small},xmin=3.5,xmax=23.5,ymin=0,ymax=26,grid=major,
  grid style={gray!25},legend cell align=left,legend style={font=\footnotesize,at={(0.02,0.97)},anchor=north west,fill=white,draw=none},
  table/col sep=comma]
\addplot[omThm,very thick,mark=*,mark size=1.6pt] table[x=quarter,y=raw]{figdata/ml/15-alternative-data-pipelines/errors.csv};
\addplot[omDef,very thick,mark=square*,mark size=1.6pt] table[x=quarter,y=raked]{figdata/ml/15-alternative-data-pipelines/errors.csv};
\addplot[black!60,dashed,ycomb,forget plot,mark=none] table[x=quarter,y expr=26]{figdata/ml/15-alternative-data-pipelines/events.csv};
\legend{spend per panelist,raked to the population margins}
\end{axis}
\end{tikzpicture}

\omcaption{Mean absolute error, across 100 retailers, of the year-on-year sales growth read from the panel. Dashed: the
vendor's launch (quarter 12) and its change of bank partner (quarter 16). Data: \texttt{ml\_altdata.errors}.}\label{fig:ml:alt:errors}
\end{omfigure}

Raking needs the population margins to be known and the panel to cover every cell: a cell the panel does not reach cannot
be reweighted into existence. It also assumes that within a cell the panel spends like the population, which is the
assumption a new partner is most likely to break (its cardholders may be richer as well as older). The vendor's own
reweighting, when there is one, deserves the same scrutiny as the raw panel: ask for the margins it rakes to and for the
panel's composition by quarter.

\section{Detecting backfill}

\begin{definition}[First-seen timestamp]\label{def:ml:alternative-data-pipelines:firstseen}
The \emph{first-seen timestamp}\index{first-seen timestamp} of a value is the time the firm first received it, recorded by
the firm at ingestion, as opposed to the valid time the vendor attaches to it. A value whose first-seen time is much later
than its valid time was backfilled (Book~7, chapter~3), and its history cannot be trusted to be what a live user would have
seen.
\end{definition}

The pipeline stores every value in the bitemporal store of Book~7 (chapter~3), with the delivery quarter as knowledge time,
and the backfill test is one query: quarters 0 to 10 were first seen at quarter 12, more than a quarter after they ended.
They matter because the vendor built them at launch with hindsight, moving its panel estimates three quarters of the way
towards the sales the retailers had already reported. (Quarter 11, delivered at launch one quarter after it ended, looks
like a live quarter and was not adjusted: its sales had not yet been reported.)

To score the nowcast as a predictor, each retailer's announcement return is its sales surprise, true growth minus a
consensus that misses by 3 points on average, plus noise. The rank information coefficient between the raked nowcast's
surprise and the return is 0.480 on the backfilled quarters and 0.419 on the live quarters before the partner change, 0.412
after it (\cref{fig:ml:alt:ic}). The raw nowcast's IC falls from 0.391 live before the change to 0.264 after it. A trial
that scores the backfilled history, which is most of what a new vendor can offer, sees a signal about 15\% stronger than
the one it will trade.

\begin{omfigure}
\begin{tikzpicture}
\begin{axis}[width=10cm,height=5.8cm,ybar,bar width=12pt,ymin=0,ymax=0.55,ylabel={mean rank IC},area legend,
  xtick={0,1,2},xticklabels={backfilled,live before\\the change,live after\\the change},tick label style={font=\small},
  xticklabel style={align=center},
  label style={font=\small},enlarge x limits=0.25,grid=major,grid style={gray!25},legend columns=2,
  legend style={font=\footnotesize,at={(0.5,-0.33)},anchor=north,draw=none,fill=none},table/col sep=comma]
\addplot[fill=omThm,draw=none] table[x=i,y=raw]{figdata/ml/15-alternative-data-pipelines/ic.csv};
\addplot[fill=omDef,draw=none] table[x=i,y=raked]{figdata/ml/15-alternative-data-pipelines/ic.csv};
\legend{spend per panelist,raked}
\end{axis}
\end{tikzpicture}

\omcaption{Mean rank information coefficient of the nowcast's surprise with the announcement return, by period. Data:
\texttt{ml\_altdata.ic\_summary}.}\label{fig:ml:alt:ic}
\end{omfigure}

\section{From pipeline to predictor}

The nowcast becomes a predictor through the vendor-trial harness of Book~7 (chapter~12): coverage by quarter, the backfill
share, a panel-drift test on the residuals against reported sales, the IC by year and the incremental IC over what the desk
already owns. The pipeline's job is to hand it data that mean what they say. Froot, Kang, Ozik and Sadka (2017) built
real-time sales proxies for retailers from about 50 million mobile devices and related them to earnings surprises and
announcement returns; Katona, Painter, Patatoukas and Zeng (2024) used the introduction of satellite coverage of retailers'
car parks to show that investors with such data traded profitably ahead of retailers' reports, especially bad ones. Both
depend on the plumbing this chapter builds: which stores belong to which company, who is in the panel, and when each number
was known.

\begin{method}[An alternative-data pipeline]\label{met:ml:alt:method}
\begin{enumerate}
\item Validate every row against a schema and every delivery against the last; reject with reasons, hold what looks wrong,
never repair silently.
\item Resolve entities to permanent identifiers with a match threshold, a review band and a person on the queue; measure
precision on a labelled sample.
\item Store every value with its first-seen time; label backfill; score trials on live data only.
\item Obtain the panel's composition by period and rake it to known margins; monitor composition for breaks.
\item Hand the result to the trial harness and, once live, to monitoring (chapter~27).
\end{enumerate}
\end{method}

\section{Tutorial: the panel that grew}\label{tut:ml:alternative-data-pipelines:tut}

\begin{tutorial}
\textbf{Goal.} Ingest a vendor's deliveries, resolve its merchant strings, rake its panel and compare nowcast errors and ICs
by period. \textbf{End state:} \cref{fig:ml:alt:errors,fig:ml:alt:ic}.
\begin{tutsteps}
\item \textbf{\omterm{def:ml:alternative-data-pipelines:ingest}{Schema validation}.}
\omcode{code/firm/altdata/firm_altdata.py}{37}{60}{Row validation against a schema.}\label{lst:ml:alt:validate}
\item \textbf{Merchant strings to identifiers.}
\omcode{code/firm/altdata/firm_altdata.py}{77}{112}{Token similarity and resolution with a review band.}\label{lst:ml:alt:resolve}
\item \textbf{Raking.}
\omcode{code/firm/altdata/firm_altdata.py}{115}{126}{Iterative proportional fitting to two margins.}\label{lst:ml:alt:rake}
\item \textbf{Run} \texttt{ml\_altdata.ingest()}, \texttt{resolution\_quality()}, \texttt{error\_summary()},
\texttt{ic\_summary()} and \texttt{fig\_altdata.py}.
\end{tutsteps}
\textbf{What to change next.} Give the new partner's cardholders a different spending level within each cell, and watch
raking fail; lower the review threshold and count the wrong matches.
\end{tutorial}

\section{Build: the alternative-data pipeline}\label{bld:ml:alternative-data-pipelines:altdata}

\begin{build}
\bfield{Purpose} Vendor deliveries turned into point-in-time, identifier-keyed, reweighted data a predictor can use.
\bfield{Interface} \texttt{SCHEMA}, \texttt{validate(rows, schema)}, \texttt{delivery\_check(prev, new, lo, hi)};
\texttt{normalise}, \texttt{similarity}, \texttt{resolve(strings, aliases, accept, review)}; \texttt{rake(counts, row\_margin,
col\_margin, iters)}; \texttt{first\_seen(store, entity, field, valid)} on \texttt{firm.pit}; \texttt{card\_panel(n\_firms,
quarters, launch, partner\_change, seed)}.
\bfield{Rules} Rejections carry reasons; nothing is repaired without a record; every stored value has a knowledge time;
matches below the threshold wait for a person.
\bfield{Acceptance tests} \texttt{code/firm/altdata/tests/}: each defect type is rejected with its reason and clean rows pass;
a doubled delivery is held; truncated and prefixed strings resolve, shared stems go to review, unrelated merchants do not
match; raking matches both margins and leaves a panel already at the margins unweighted; first-seen times separate a
backfilled history.
\bfield{Stretch} Fellegi--Sunter match weights learned from the review decisions; raking to three margins with a cap on
weights; the panel's composition monitored for breaks.
\end{build}

\begin{omsources}
\item I.~P.~Fellegi and A.~B.~Sunter, ``A theory for record linkage'', \emph{Journal of the American Statistical Association}
64(328), 1969.
\item W.~E.~Deming and F.~F.~Stephan, ``On a least squares adjustment of a sampled frequency table when the expected marginal
totals are known'', \emph{Annals of Mathematical Statistics} 11(4), 1940.
\item K.~A.~Froot, N.~Kang, G.~Ozik and R.~Sadka, ``What do measures of real-time corporate sales say about earnings
surprises and post-announcement returns?'', \emph{Journal of Financial Economics} 125(1), 2017.
\item Z.~Katona, M.~Painter, P.~N.~Patatoukas and J.~Zeng, ``On the capital market consequences of big data: evidence from
outer space'', \emph{Journal of Financial and Quantitative Analysis}, 2024.
\end{omsources}

\section{Exercises}

\begin{exercise}[$\star$]\label{exo:ml:alternative-data-pipelines:1}
A delivery totals 2.4 billion; the previous one totalled 23.9 million. What does the delivery check do, and what are the
likely causes?
\end{exercise}

\begin{exercise}[$\star$]\label{exo:ml:alternative-data-pipelines:2}
Rake the $2\times2$ panel counts $\begin{pmatrix}30 & 20\\ 10 & 40\end{pmatrix}$ to even row and column margins. Give the
weighted counts.
\end{exercise}

\begin{exercise}[$\star$]\label{exo:ml:alternative-data-pipelines:3}
Why does ``BOREALIS 5120'' go to review while ``BOREALI RETA STR 313'' is matched?
\end{exercise}

\begin{exercise}[$\star\star$]\label{exo:ml:alternative-data-pipelines:4}
Explain, with the chapter's world, why the raw nowcast's error rises for exactly four quarters after the partner change.
\end{exercise}

\begin{exercise}[$\star\star$]\label{exo:ml:alternative-data-pipelines:5}
Why could the first-seen test not flag quarter 11, and why did that not matter here? When would it?
\end{exercise}

\begin{exercise}[$\star\star$]\label{exo:ml:alternative-data-pipelines:6}
\emph{Find the flaw.} ``We fixed the vendor's negative spends by taking absolute values and cast the text panel counts to
integers, so no data were lost.''
\end{exercise}

\begin{exercise}[$\star\star\star$]\label{exo:ml:alternative-data-pipelines:7}
\emph{Coding.} Run the pipeline with the review queue left unresolved (strings in review dropped). What happens to the raked
nowcast's error, and why is the error larger for some retailers than others?
\end{exercise}

\begin{exercise}[$\star\star\star$]\label{exo:ml:alternative-data-pipelines:8}
Show that with a panel whose cell counts are the product of an age margin and a region margin, raking to the population's
margins gives the post-stratification weights (population share over panel share of each cell) when the population's cell
shares are also a product. What happens when they are not?
\end{exercise}

\section{Problem: The Panel That Grew}

\begin{problem}[{Weekend problem --- a panel is not a population}]\label{pb:ml:alternative-data-pipelines:1}
The chapter's world: 100 retailers, 24 quarters, a panel that drifts, changes partner at quarter 16 and launches at quarter
12 with a backfilled history.

\medskip\noindent\textbf{Part I --- Ingestion.}
\begin{enumerate}
\item What defects were planted, and how many rows does validation reject?
\item Which defect does row validation miss, and what catches it?
\item Why reject rather than repair?
\item What does the pipeline store for each value, and why?
\end{enumerate}

\medskip\noindent\textbf{Part II --- Resolution.}
\begin{enumerate}[resume]
\item How many strings does exact matching resolve, and how many the resolver?
\item What goes to review, and why is that the right place for it?
\item What happens to the unrelated merchants?
\item What does a wrong match cost, compared with a string left in review?
\end{enumerate}

\medskip\noindent\textbf{Part III --- Panel and backfill.}
\begin{enumerate}[resume]
\item What are the raw and raked nowcast errors before, during and after the change year?
\item What does raking assume, and what residual error remains?
\item Which quarters are backfilled, and how does the pipeline know?
\item What are the ICs on backfilled and live history?
\end{enumerate}

\medskip\noindent\textbf{Part IV --- The verdict.}
\begin{enumerate}[resume]
\item State the \emph{named result}: the nowcast error with and without reweighting after the partner change, and the IC of
the backfilled history against the live one.
\item What would you ask a card-panel vendor for before a trial?
\item How would you detect a partner change the vendor did not announce?
\item What would you monitor once the data are live?
\item How much of the backfilled IC would you believe?
\item Where would a person sit in this pipeline, and why there?
\item What changes for a satellite or app-download panel?
\item In one sentence: what is the most expensive assumption in alternative data?
\end{enumerate}
\end{problem}

\section{Interview questions}

\begin{interviewq}[$\star$ \iqroles{researcher, mle}]\label{iq:ml:alternative-data-pipelines:1}
What is backfill in a vendor dataset, and how do you detect it?
\end{interviewq}

\begin{interviewq}[$\star\star$ \iqroles{mle, dev}]\label{iq:ml:alternative-data-pipelines:2}
Design the ingestion step for a daily vendor file. What do you check, and what do you do when a check fails?
\end{interviewq}

\begin{interviewq}[$\star\star$ \iqroles{mle}]\label{iq:ml:alternative-data-pipelines:3}
How would you map merchant descriptors to listed companies at scale, and how would you measure the mapping's quality?
\end{interviewq}

\begin{interviewq}[$\star\star$ \iqroles{researcher}]\label{iq:ml:alternative-data-pipelines:4}
A card panel over-represents young customers. How do you correct a sales nowcast for it, and what can go wrong?
\end{interviewq}

\begin{interviewq}[$\star\star$ \iqroles{researcher, trader}]\label{iq:ml:alternative-data-pipelines:5}
A nowcast that tracked sales for three years misses badly for a year. What do you check first?
\end{interviewq}

\begin{interviewq}[$\star\star\star$ \iqroles{researcher}]\label{iq:ml:alternative-data-pipelines:6}
A vendor's trial shows an IC of 0.5 over eight years; the dataset launched two years ago. How would you estimate the IC you
should expect live?
\end{interviewq}
